% Math in a box saved with \sbox, \savebox, or lrbox, and placed in a TikZ
% node. Each case's text is "Before <case>." and "After <case>.", which
% run.mjs reads.
%
% Pictures whose math is drawn in the picture, as in the PDF:
% - \sbox and \savebox holding inline math;
% - an environment that saves a minipage with lrbox and draws it in a
%   node, holding inline math with $ and \(, a \[ \] display, a $$
%   display with the same formula, an equation and an align* with text
%   after them, a flalign, a display with an alignment and inline math
%   with a matrix, both with &, and an enumerate;
% - a saved minipage holding a tikzcd display, drawn once;
% - \sbox holding a siunitx quantity;
% - math typed directly in a node.
%
% Copies of saved boxes outside any picture, whose math is the page's
% formulas: in text, in \fbox, a minipage with \[ \], $$, and a labelled
% equation, which a \eqref outside the box cites, a minipage with the
% tikzcd display, and a proof ending in a display with \qedhere.
\documentclass{article}
\usepackage{amsmath}
\usepackage{amsthm}
\usepackage{tikz}
\usepackage{tikz-cd}
\usepackage{siunitx}
\newsavebox{\bxa}
\newsavebox{\bxb}
\newsavebox{\bxc}
\newsavebox{\bxd}
\newsavebox{\bxe}
\newsavebox{\bxf}
\newsavebox{\mathbx}
\newenvironment{mathbox}{\begin{lrbox}{\mathbx}\begin{minipage}{0.8\textwidth}}{%
  \end{minipage}\end{lrbox}\begin{center}\begin{tikzpicture}\node[draw, inner sep=10pt]{\usebox{\mathbx}};\end{tikzpicture}\end{center}}
\begin{document}

Before sbox. \sbox{\bxa}{Box $x^2$ end}\tikz\node[draw]{\usebox{\bxa}}; After sbox.

Before savebox. \savebox{\bxb}[4cm][l]{Saved \(y_1 + z\) here}\tikz\node[draw]{\usebox{\bxb}}; After savebox.

Before lrbox inline.
\begin{mathbox}
Let $p$ be a prime and \(q = p^2\) its square, so that every line of this
paragraph breaks where the formulas' widths put it.
\end{mathbox}
After lrbox inline.

Before lrbox display.
\begin{mathbox}
The sum is
\[ \sum_{k=1}^{n} k = \frac{n(n+1)}{2}, \]
for every whole number $n$.
\end{mathbox}
After lrbox display.

Before lrbox dollars.
\begin{mathbox}
The sum is
$$ \sum_{k=1}^{n} k = \frac{n(n+1)}{2}, $$
for every whole number $n$.
\end{mathbox}
After lrbox dollars.

Before lrbox equation.
\begin{mathbox}
A numbered formula:
\begin{equation}
  e^{i\pi} + 1 = 0.
\end{equation}
which joins five constants.
\end{mathbox}
After lrbox equation.

Before lrbox align.
\begin{mathbox}
Two rows:
\begin{align*}
  a &= b + c, \\
  d &= e - f.
\end{align*}
which hold together.
\end{mathbox}
After lrbox align.

Before lrbox flalign.
\begin{mathbox}
Flush rows:
\begin{flalign} a &= b && \end{flalign}
\end{mathbox}
After lrbox flalign.

Before lrbox ampersand.
\begin{mathbox}
A matrix $\left(\begin{smallmatrix} u & v \\ w & t \end{smallmatrix}\right)$ and
\[ \begin{aligned} g &= h, \\ m &= n. \end{aligned} \]
\end{mathbox}
After lrbox ampersand.

Before lrbox list.
\begin{mathbox}
\begin{enumerate}
\item The first item, with $r < s$.
\item The second item, with \(s < t\).
\end{enumerate}
\end{mathbox}
After lrbox list.

\begin{lrbox}{\bxd}\begin{minipage}{0.5\textwidth}
arrows:
\[ \begin{tikzcd} A \arrow[r] & B \end{tikzcd} \]
\end{minipage}\end{lrbox}
Before lrbox tikzcd. \tikz\node[draw]{\usebox{\bxd}}; After lrbox tikzcd.

Before sbox quantity. \sbox{\bxf}{Length \qty{3}{\metre} end}\tikz\node[draw]{\usebox{\bxf}}; After sbox quantity.

Before direct. \tikz\node[draw]{Node $y^2$ end}; After direct.

Before text copy. Kept \usebox{\bxa} here. After text copy.

Before fbox copy. Framed \fbox{\usebox{\bxb}} here. After fbox copy.

\begin{lrbox}{\bxc}\begin{minipage}{0.5\textwidth}
Copied $v$ and
\[ w = 2v, \]
then
$$ u = 3v, $$
and
\begin{equation}\label{eq:boxed} t = 4v. \end{equation}
\end{minipage}\end{lrbox}
Before minipage copy. Placed \usebox{\bxc} here. After minipage copy.

Before tikzcd copy. Placed \usebox{\bxd} here. After tikzcd copy.

\begin{lrbox}{\bxe}\begin{minipage}{0.5\textwidth}
\begin{proof}
Then
\[ x = y. \qedhere \]
\end{proof}
\end{minipage}\end{lrbox}
Before qed copy. Placed \usebox{\bxe} here. After qed copy.

Before eqref. See \eqref{eq:boxed}. After eqref.

\end{document}
